% Chapter 5: Clock Management
% File: chapters/chapter-05-clocks/chapter-05-clocks.tex

\chapter{Phase A: Clock Management and System Configuration}
\label{ch:chapter-05}
\chapthumb{5}

\begin{quote}
    \textit{``Every timing bug is a clock bug wearing a disguise.''}
\end{quote}

Nothing on this part runs without a clock, and almost nothing that goes wrong is blamed on one. This
chapter is the record of four separate defects in one project that were all clocks: a delay that was
three times too short because a loop is not an instruction count, a peripheral that never started
because its oscillator was selected but not enabled, a converter that hung because it was clocked
from the wrong source, and a sampling loop that claimed 400 Hz and ran at 158.

The numbers here are measured with the cycle counter, not read from a datasheet.

\minitoc

\clearpage

\section{The Tree, in the Order You Meet It}
\label{sec:clocks-tree}

A Series 2 clock tree is not complicated once it is drawn as three questions. This section answers
them for this board only, because the full diagram in the reference manual has forty boxes and no
project needs forty.

\subsection{Question 1: Which Oscillator}

\begin{table}[htbp]
    \centering
    \footnotesize
    \begin{tabular}{@{}p{0.20\textwidth}p{0.16\textwidth}p{0.52\textwidth}@{}}
        \toprule
        \textbf{Oscillator} & \textbf{Rate} & \textbf{When it is the right answer} \\
        \midrule
        \codeinline{HFRCO0} & 19 MHz default & the reset state; enough for GPIO, timers and I2C \\
        \codeinline{HFXO0}  & 38.4 MHz crystal & radio, USB, anything needing ppm accuracy \\
        \codeinline{FSRCO}  & 20 MHz fast RC & cheap and fast, but see the IADC below \\
        \codeinline{HFRCOEM23} & 19 MHz & low-energy domain; the IADC and TIMER0 like it \\
        \codeinline{DPLL0}  & up to 78 MHz & the CPU's fast mode, locked to HFRCO or HFXO \\
        \codeinline{LFRCO}/\codeinline{LFXO} & 32.768 kHz & RTC, sleep, wake-up \\
        \bottomrule
    \end{tabular}
    \caption{Six oscillators, and the four that matter for a board with an I2C bus, a buzzer and a
    microphone.}
    \label{tab:clocks-oscillators}
\end{table}

\subsection{Question 2: Which Branch}

Each oscillator feeds branches, and each branch has a selector and a divider:

\begin{minted}{rust}
// The low-energy peripherals (IADC, TIMER0 in this project) are clocked from EM01GRPACLK.
dp.cmu_s.em01grpaclkctrl().modify(|_r, w| w.clksel().hfrcoem23());
\end{minted}

The name matters when reading the manual: \codeinline{EM01GRPACLK} is one branch,
\codeinline{EM23GRPACLK} another, and a peripheral belongs to exactly one. Choosing the right
oscillator on the wrong branch gives a peripheral clocked at zero — which looks identical to a
peripheral whose registers are not being written.

\subsection{Question 3: Which Gate}

Every peripheral has a bus-clock gate in \codeinline{CMU\_CLKEN0} or \codeinline{CMU\_CLKEN1}. On
this part the GPIO bank is bit 0 of \codeinline{CLKEN0} and the MVP's is bit 30 of
\codeinline{CLKEN1}. A register access with the gate closed is a bus fault, precisely reported:

\begin{minted}{text}
CFSR = 0x00008200   # PRECISERR set: a precise data bus fault
BFAR = 0x5003C040   # ...at PORTA_DOUT, whose clock was off
\end{minted}

That pair of registers is the fastest diagnosis in this book. If a peripheral reads as zero, check
the gate; if a write faults, read \codeinline{CFSR} and \codeinline{BFAR} before anything else.

\subsection{The Three Questions, Together}

\begin{minted}{rust}
// Oscillator: force HFRCOEM23 on and wait for it.
dp.cmu_s.clken0().modify(|_r, w| w.hfrcoem23().set_bit());
dp.hfrcoem23_s.ctrl().modify(|_r, w| w.forceen().set_bit());
while !dp.hfrcoem23_s.status().read().rdy().bit_is_set() {}

// Branch: the low-energy domain comes from it.
dp.cmu_s.em01grpaclkctrl().modify(|_r, w| w.clksel().hfrcoem23());

// Gate: the peripheral's own bus clock.
dp.cmu_s.clken0().modify(|_r, w| w.timer0().set_bit());
\end{minted}

Three decisions, three writes, and one wait. The wait is the subject of the next-but-one section: it
is the only one of the four steps that can hang.

\section{The Clock You Forgot You Were Using}
\label{sec:clocks-sysclk}

Every HAL has a constant called something like \codeinline{DEFAULT\_SYSCLK\_HZ}, and it is almost
always wrong for the board it runs on. This section is about the one on this board, because it
silently defines the length of every delay in the project.

\subsection{19 MHz, Not 78}

The EFR32MG24 can run at 78 MHz. It does not, out of reset, and this board's firmware never raises
it. The default is \codeinline{HFRCO0} at 19 MHz, and the measurement that settled it was a stopwatch
rather than a datasheet:

\begin{minted}{text}
# a blink with a 500 ms half-period, counted against a phone stopwatch
# 20 blinks took 20.0 s  ->  1.000 s per cycle  ->  the delay was right at 19 MHz
# the same code timed as if SYSCLK were 78 MHz produced a 4x fast blink
\end{minted}

That is a crude instrument, and it is the right one: a stopwatch over twenty cycles has an error far
below the 4× discrepancy being tested.

\subsection{Why the Error Is Invisible}

A wrong \codeinline{SYSCLK\_HZ} does not produce a wrong-looking program. It produces a program whose
delays are all scaled by the same factor — blinks that are too fast, tones that are too high, buses
that are too slow or too fast. Each of those looks like a bug in that peripheral:

\begin{table}[htbp]
    \centering
    \footnotesize
    \begin{tabular}{@{}p{0.30\textwidth}p{0.30\textwidth}p{0.32\textwidth}@{}}
        \toprule
        \textbf{Symptom} & \textbf{Blamed on} & \textbf{Actually} \\
        \midrule
        LED blinks at 4× the rate & the delay loop & \codeinline{SYSCLK\_HZ} too high \\
        Tone is 4× the intended pitch & the timer's TOP & same constant \\
        I2C bus works on one board, not another & the devices & same constant, other clock \\
        Buzzer sounds, then the beep pattern is wrong & the pattern code & same constant \\
        \bottomrule
    \end{tabular}
    \caption{One wrong constant, four different suspects. This is why it is measured once and written
    down next to the definition.}
    \label{tab:sysclk-symptoms}
\end{table}

\subsection{The Right Shape for the Constant}

The version that survives contact with a second board is not a number in the HAL. It is a number the
board passes in:

\begin{minted}{rust}
// The HAL knows how to delay; it does not know how fast this board's clock runs.
let mut delay = Delay::with_sysclk(syst, 19_000_000);
\end{minted}

and in this project's board layer, one place says it out loud:

\begin{minted}{rust}
pub const DEFAULT_SYSCLK_HZ: u32 = 19_000_000;   // measured, see 05.02
\end{minted}

A second board with a 38.4 MHz crystal changes one line, and every delay, tone and bus timing
follows it. A HAL with the constant baked in cannot be reused, and — worse — fails in a way that
looks like the new board's fault.

\subsection{Measuring It Properly}

The stopwatch is enough to catch a factor of four. For anything finer, use the cycle counter, which
counts CPU cycles and therefore cannot disagree with the CPU:

\begin{minted}{rust}
let start = dwt::cycles();
delay.delay_ms(1000);
let elapsed = dwt::since(start);
// elapsed / 19_000_000 must be ~1.0 s; a factor of 3.0 here means the delay's
// internal arithmetic is wrong, not the constant.
\end{minted}

That distinction — constant wrong \emph{versus} arithmetic wrong — is the difference between
recalibrating and debugging, and the cycle counter is what separates them.

\section{Selecting an Oscillator Is Not Starting It}
\label{sec:clocks-starting}

The clock tree has a trap that costs an afternoon the first time and ten minutes every time after:
choosing a source does not turn it on, and a peripheral clocked by a stopped source behaves like a
peripheral that ignores you.

\subsection{The Two-Step That Is Actually Three}

\begin{minted}{rust}
// 1. Ask for it.
dp.cmu_s.clken0().modify(|_r, w| w.hfrcoem23().set_bit());
// 2. Force it on: without this the oscillator idles when nothing demands it,
//    and "nothing" includes the branch you are about to select.
dp.hfrcoem23_s.ctrl().modify(|_r, w| w.forceen().set_bit());
// 3. Wait. This is the step that hangs if the oscillator never comes up.
while !dp.hfrcoem23_s.status().read().rdy().bit_is_set() { /* spin */ }
\end{minted}

Step 2 is the one people miss. It is also the one that makes step 3 a real wait instead of an
instant pass — and an instant pass is worse than a hang, because it looks like success.

\subsection{The Converter That Hung on FSRCO}

The microphone's converter (IADC0) had two separate problems in this project, and one of them was
purely a clock choice. Clocked from \codeinline{FSRCO}, the converter's disable sequence never
completed:

\begin{minted}{text}
# after writing EN = 0 and waiting for the disabling bit to clear
IADC0_STATUS  = ...DISABLING = 1     (still, after 100 ms)
# the next configuration write then faulted, because the peripheral was mid-synchronisation
\end{minted}

The same code, clocked from \codeinline{HFRCOEM23}, disabled in microseconds and never faulted. No
register value changed; only \codeinline{em01grpaclkctrl}. That is the entire difference between a
working microphone driver and an afternoon of bus faults, and it was found by trying the other
oscillator rather than by reading harder.

\subsection{The Rule That Came Out of It}

Before blaming a peripheral, confirm that it is \emph{clocked}, and confirm it in the way the
hardware reports it:

\begin{enumerate}
    \item read the gate bit back — not the value you wrote, the value the register holds;
    \item read the oscillator's ready status, not just its enable bit;
    \item if the peripheral has a synchronisation status, read it: \codeinline{DISABLING},
          \codeinline{SYNCBUSY} and \codeinline{RESETTING} are the hardware telling you the
          difference between ``configured'' and ``ready''.
\end{enumerate}

A peripheral that reads as zero after a correct-looking write has failed one of those three, and the
three have different fixes. Without the read-backs they are indistinguishable — which is why every
driver in this repository returns an error instead of a value when a readiness wait expires.

\section{Delay Without Lying}
\label{sec:clocks-delay}

There are three ways to wait on a Cortex-M and this project has been wrong about all three.

\subsection{asm::delay Is Not a Cycle Count}

\begin{minted}{rust}
// "1000 cycles", according to the name and to every example on the internet.
cortex_m::asm::delay(1000);
\end{minted}

Measured against the cycle counter, one iteration of the loop that \codeinline{asm::delay} expands
to takes about \textbf{three} cycles on this part with this compiler. So
\codeinline{delay(1\_000\_000)} waits three million cycles, not one — a 3× error, and the reason a
2.7 kHz tone computed with it came out as something else entirely.

Two consequences, both practical:

\begin{itemize}
    \item \textbf{Never derive a period from it.} A tone, a baud rate or a bus half-period computed
          with \codeinline{asm::delay} is wrong by a factor nobody wrote down.
    \item \textbf{Do not multiply by three and call it fixed.} The factor is a property of the
          compiler's code generation, not of the hardware; it changed between optimisation levels in
          this project.
\end{itemize}

\subsection{SysTick: Real Milliseconds}

SysTick counts CPU cycles and can therefore produce real time:

\begin{minted}{rust}
pub fn delay_us(&mut self, us: u32) {
    let cycles = (u64::from(us) * u64::from(self.sysclk_hz)) / 1_000_000;
    // SysTick's reload is 24 bits: split anything larger into chunks.
    let mut remaining = cycles;
    while remaining > 0 {
        let chunk = remaining.min(0x00FF_FFFF);
        self.syst.set_reload(chunk as u32);
        self.syst.clear_current();
        self.syst.enable_counter();
        while !self.syst.has_wrapped() {}
        self.syst.disable_counter();
        remaining -= chunk;
    }
}
\end{minted}

It has one owner — the board — and it must never be shared with a scheduler that also wants it. In
this project that rule is enforced by the type system in the smallest possible way: \codeinline{SYST}
is taken once by \codeinline{Board::take()}, and \codeinline{board.delay()} hands it out.

\subsection{Ever the Millisecond: Do Not Divide Early}

A defect worth its own paragraph, because it looks like arithmetic hygiene:

\begin{minted}{rust}
// The sample interval arrives as microseconds: 2500 us, i.e. 416 Hz.
let dt_ms = dt_us / 1000;        // 2 ms
// Every timeout in the state machine is now 20% short.
\end{minted}

Integer division by 1000 truncated 2500 to 2, and every threshold derived from it inherited a 20%
error. The classifier's free-fall window, its landing window and its settle timer were all wrong at
once, and every one of them was measured against real drops before anyone suspected the unit.

The rule: \textbf{carry the finest unit the hardware gives you and divide only where a human reads
it}. This project's classifier now takes microseconds everywhere and converts once, at the boundary,
for display.

\subsection{The Cycle Counter: Measurement, Not Delay}

The third instrument is not for waiting at all:

\begin{minted}{rust}
let start = dwt::cycles();
// ... work ...
let elapsed_us = dwt::cycles_to_us(dwt::since(start), 19_000_000);
\end{minted}

Two registers and no ownership questions: enable \codeinline{TRCENA} in \codeinline{DEMCR}, then
\codeinline{CYCCNTENA} in \codeinline{DWT->CTRL}, then read \codeinline{CYCCNT}. It is how every
timing number in this book was obtained, and it is the only instrument here that cannot lie about
the clock because the CPU \emph{is} the clock.

One trap, measured: a counter that was never enabled reads as zero, and zero elapsed cycles looks
exactly like a loop that ran impossibly fast. Compare two consecutive reads before trusting either.

\section{What to Write Down, and Where}
\label{sec:clocks-write-down}

A clock fact has three places it can live, and only one of them survives the project.

\subsection{The Three Places}

\begin{table}[htbp]
    \centering
    \footnotesize
    \begin{tabular}{@{}p{0.28\textwidth}p{0.30\textwidth}p{0.34\textwidth}@{}}
        \toprule
        \textbf{Place} & \textbf{Lifetime} & \textbf{What happens to it} \\
        \midrule
        In the chat where you found it & minutes & lost \\
        In a commit message & months, if you search & nobody searches \\
        In the source, next to the constant it explains & the life of the code & \textbf{read by the next person} \\
        \bottomrule
    \end{tabular}
    \caption{The only durable location is the one somebody will read while editing the thing it
    describes.}
    \label{tab:where-facts-live}
\end{table}

\subsection{How the Numbers Here Are Written}

Every measured constant in this project carries three things: the number, the unit, and how it was
obtained. For example:

\begin{minted}{rust}
/// Half a bit period on the bit-banged I2C bus, in loop iterations.
///
/// Measured: 60 iterations take ~5 us at SYSCLK = 19 MHz on this board, so ~95 kHz.
/// This is a measurement, not a datasheet value -- `asm::delay` is about three cycles
/// per iteration here, so the divisor is not the clock.
const HALF_BIT_CYCLES: u32 = 60;
\end{minted}

and for the value that everything else depends on:

\begin{minted}{rust}
/// SYSCLK as measured with a stopwatch over 20 blink cycles: 19 MHz (the HFRCO default).
/// The part supports 78 MHz; this board's firmware never selects it.
pub const DEFAULT_SYSCLK_HZ: u32 = 19_000_000;
\end{minted}

\subsection{The Table That Belongs in Every Project}

\begin{table}[htbp]
    \centering
    \footnotesize
    \begin{tabular}{@{}p{0.28\textwidth}p{0.24\textwidth}p{0.40\textwidth}@{}}
        \toprule
        \textbf{Quantity} & \textbf{Value here} & \textbf{How it was measured} \\
        \midrule
        SYSCLK & 19 MHz & stopwatch over 20 blink cycles \\
        \codeinline{asm::delay} cost & ~3 cycles/iteration & cycle counter, two build profiles \\
        I2C half period & ~5 $\mu$s (60 iterations) & scope-free: 95 kHz against a known device \\
        Buzzer tone & 2.7 kHz, TOP = 3518 & TIMER0 TOP read back over SWD \\
        IMU sample rate & 416 Hz configured & \codeinline{CTRL1\_XL} register + measured intervals \\
        Sampling loop & 158 Hz when claiming 400 & cycle counter, per-sample interval recorded \\
        \bottomrule
    \end{tabular}
    \caption{The clock facts of this project. The last line is the one that mattered: a loop that
    reports a rate it does not achieve poisons every timing test downstream.}
    \label{tab:clock-facts}
\end{table}

\subsection{The One That Generalises}

The last row of that table is the lesson of this chapter. The firmware said 400 Hz because the
\emph{intended} interval was 2500 $\mu$s. It ran at 158 Hz because the display flush blocked for tens
of milliseconds inside the loop. Every threshold in the state machine was therefore 2.5× fast, and
three rounds of drop tests were interpreted against a wrong time base before anyone noticed that four
different drops showed the same 0.62 s of calm.

The fix was not a better delay. It was to \textbf{measure the interval between samples and pass the
measurement on}:

\begin{minted}{rust}
let now = dwt::cycles();
let dt_us = dwt::cycles_to_us(dwt::since(self.last), sysclk_hz);
self.last = now;
// The classifier is told what actually happened, not what was planned.
detector.update(&accel, dt_us);
\end{minted}

Every timing defect in this project — and there were four — would have been caught on the day it was
introduced by that three-line habit.


\clearpage

\begin{keytakeaways}
    \item A clock tree is three separate decisions: which oscillator, which branch, and which
          peripheral gate. Getting one of the three wrong produces the same symptom — a peripheral
          that reads as zero.
    \item SYSCLK on this board is the HFRCO default, measured at 19 MHz. The part can do 78 MHz; the
          board does not, until someone configures it.
    \item Selecting an oscillator is not starting it. Force it on, then wait for the ready bit.
    \item \codeinline{asm::delay(n)} is not \codeinline{n} cycles — it measured about three per
          iteration here, which is a 3× error in every derived period.
    \item A loop's \emph{intended} rate is not its actual rate. Measure the interval between samples
          and pass the measurement on, never the intention.
\end{keytakeaways}
